% !TeX root = ../cosmology_summary.tex

%% --------------------------------------------------------
\chapter{Збурення метрики: SVT-розклад\\ і калібрування}
\label{sec:metric}
%% --------------------------------------------------------

Рівняння Фрідмана описують ідеально симетричний фон.
Реальний Всесвіт відхиляється від нього: є флуктуації густини,
потоки речовини, гравітаційні хвилі. Щоб описати ці відхилення
кількісно~--- і зв'язати їх з CMB~--- потрібна мова малих збурень
навколо метрики FLRW. Ця мова~--- лінійна теорія збурень метрики~---
стане фундаментом для квантування збурень у наступному розділі.

Малі відхилення $\delta g_{\mu\nu}(\eta, \mathbf{x})$ від
фонового значення $\bar{g}_{\mu\nu}(\eta)$ класифікуються за
їхньою поведінкою під просторовими обертаннями. Ключовий факт:
оскільки фон є однорідним і ізотропним, збурення різних
тензорних рангів не взаємодіють між собою у лінійному порядку~---
скалярні, векторні і тензорні компоненти еволюціонують
незалежно. Це твердження відоме як SVT-розклад
\cite{stewart_walker_1974}:
\begin{equation}\label{eq:svt_decomposition}
  \delta g_{\mu\nu}
    = \delta g_{\mu\nu}^{(\mathrm{S})}
    + \delta g_{\mu\nu}^{(\mathrm{V})}
    + \delta g_{\mu\nu}^{(\mathrm{T})}.
\end{equation}
Незалежність трьох секторів означає, що кожен з них можна
вивчати окремо:
\begin{itemize}
  \item \textbf{Скалярні збурення ($\mathrm{S}$):}
        Параметризуються двома скалярними функціями
        $(\Phi, \Psi)$. Динамічно пов'язані з флуктуаціями
        густини матерії $\delta\rho$~--- це джерело
        великомасштабної структури.
  \item \textbf{Векторні збурення ($\mathrm{V}$):}
        Описують вихрові рухи і моменти імпульсу. За
        відсутності джерел вони згасають як $a^{-2}$
        \cite{mukhanov_2005} і в стандартній моделі ролі не
        відіграють. Ми вводимо їх тут для повноти
        класифікації і далі не розглядаємо.
  \item \textbf{Тензорні збурення ($\mathrm{T}$):}
        Вільні поперечні безслідові деформації просторової
        метрики~--- первинні гравітаційні хвилі. Народжуються
        під час інфляції і практично не взаємодіють з
        речовиною.
\end{itemize}

%% --------------------------------------------------------
\section{Конформно-ньютонівське калібрування}
%% --------------------------------------------------------

Рівняння загальної теорії відносності є загально-коваріантними, що забезпечує
свободу вибору просторово-часових координат. У лінеаризованій теорії
космологічних збурень ця коваріантність призводить до появи суто
калібрувальних (gauge) ступенів вільності, які є артефактами координатного
опису і позбавлені безпосереднього фізичного змісту. Для усунення цієї
неоднозначності та фіксації фізичних мод ми використовуємо
\emph{конформно-ньютонівське} (поздовжнє) калібрування. У цьому формулюванні
позадіагональні компоненти метрики тотожно дорівнюють нулю, а збурений
інтервал набуває вигляду:
\begin{equation}\label{eq:perturbed_metric}
  ds^2 = a^2(\eta)\Bigl[
    -(1 + 2\Phi)\,d\eta^2
    + \bigl((1 - 2\Psi)\,\delta_{ij} + h_{ij}\bigr)\,dx^i dx^j
  \Bigr].
\end{equation}
Тут скалярні флуктуації геометрії описуються двома потенціалами
Бардіна $\Phi$ та $\Psi$ \cite{bardeen_1980}, які у квазістатичному
граничному випадку на масштабах під горизонтом зводяться до класичного
ньютонівського потенціалу, а $h_{ij}$ відповідає за тензорні збурення.

Фізичний зміст кожного доданку:
\begin{enumerate}
  \item $\Phi$~--- ньютонівський гравітаційний потенціал:
        визначає прискорення пробних частинок і гравітаційне
        червоне зміщення фотонів CMB.
  \item $\Psi$~--- потенціал просторової кривини: описує
        локальне стиснення або розтягнення просторових
        гіперповерхонь. За відсутності анізотропного напруження
        $\Phi = \Psi$.
  \item $h_{ij}$~--- задовольняє умовам поперечності
        ($\partial^i h_{ij} = 0$) і безслідовості
        ($h^i{}_i = 0$); описує два стани поляризації
        гравітаційних хвиль $h_+$ та $h_\times$.
\end{enumerate}

Таким чином, цей розділ побудував мову: збурена метрика
параметризована потенціалами $\Phi$, $\Psi$ і тензором $h_{ij}$,
а SVT-розклад гарантує їхню незалежну еволюцію. Але мова збурень
метрики корисна лише тоді, коли ми знаємо, як ця збурена геометрія
впливає на речовину. Покажемо це явно~--- отриманий результат стане
фундаментом для гідродинаміки фотон-баріонної плазми у
розділі~\ref{sec:bao}.

%% --------------------------------------------------------

Тепер можна поставити центральне питання: звідки взялися
самі збурення $\delta\rho$, $\Phi$~--- початкові умови
для цієї гідродинаміки? Відповідь~--- квантові флуктуації
інфлатона~--- є темою наступного розділу.



% ============================================================
% Задачі до розділу 3: Збурення метрики
\section*{Задачі}
\addcontentsline{toc}{section}{Задачі}
% ============================================================

\begin{problem}{SVT-розклад та незалежність секторів}
У лінійній теорії космологічних збурень метрику FLRW з малими відхиленнями
розкладають на скалярний, векторний та тензорний сектори.

\begin{enumerate}
    \item Скільки незалежних ступенів вільності має кожен із секторів
          (S, V, T) у збуреній метриці? Поясніть, чому тензорні збурення
          описують лише дві поляризації гравітаційних хвиль.
    \item Чому скалярні, векторні та тензорні збурення еволюціонують
          незалежно в лінійному порядку? Яка властивість фонового
          простору-часу за це відповідає?
    \item Векторні збурення за відсутності джерел згасають як $a^{-2}$.
          Поясніть фізичну причину цього згасання. Чому в стандартній
          космології вони зазвичай не розглядаються?
\end{enumerate}
\end{problem}

% -------------------------------------------------------------------

\begin{problem}{Конформно-ньютонівське калібрування та потенціали Бардіна}
У конформно-ньютонівському (поздовжньому) калібруванні збурена метрика має вигляд:
\[
ds^2 = a^2(\eta)\left[-(1 + 2\Phi)\,d\eta^2 + \bigl((1 - 2\Psi)\,\delta_{ij} + h_{ij}\bigr)\,dx^i dx^j\right].
\]

\begin{enumerate}
    \item Поясніть фізичний зміст потенціалів $\Phi$ та $\Psi$.
          У якому випадку вони збігаються ($\Phi = \Psi$)?
    \item Що таке "калібрувальна свобода" в загальній теорії відносності?
          Чому для усунення неоднозначності потрібно фіксувати калібрування?
    \item Уявіть, що ви обрали інше калібрування (наприклад, синхронне).
          Чи будуть фізичні передбачення (наприклад, кутовий спектр CMB)
          залежати від вибору калібрування? Відповідь обґрунтуйте.
\end{enumerate}
\end{problem}

\begin{problem}{Тензорні збурення та гравітаційні хвилі}
Тензорні збурення $h_{ij}$ задовольняють умовам поперечності ($\partial^i h_{ij}=0$)
та безслідовості ($h^i{}_i=0$).

\begin{enumerate}
    \item Скільки незалежних компонент залишається в $h_{ij}$ після накладання
          цих умов у тривимірному просторі?
    \item Чому ці незалежні компоненти відповідають двом поляризаціям
          гравітаційних хвиль ($h_+$ та $h_\times$)?
    \item Покажіть, що для плоскої хвилі, яка поширюється вздовж осі $z$,
          тензор $h_{ij}$ має вигляд:
          \[
          h_{ij} = \begin{pmatrix} h_+ & h_\times & 0 \\ h_\times & -h_+ & 0 \\ 0 & 0 & 0 \end{pmatrix}.
          \]
          Як трансформуються $h_+$ та $h_\times$ при повороті на кут $45^\circ$ навколо осі $z$?
\end{enumerate}
\end{problem}